% Etapa 4D: añadir las figuras a la carpeta figuras/ de Overleaf.
% La evaluación no constituye una verificación estructural final de la matriz.
\subsection{Simulación estructural preliminar de la matriz perforada}

Se desarrolló un modelo tridimensional de elasticidad lineal por elementos finitos
para estudiar la respuesta inicial de la matriz frente a las cargas de compactación.
La geometría representada corresponde a un disco de $200\,\mathrm{mm}$ de diámetro y
$36\,\mathrm{mm}$ de espesor, con alojamiento central de $50\,\mathrm{mm}$ y
$102$ perforaciones pasantes de $6\,\mathrm{mm}$ posicionadas según los centros del CAD.
Se consideraron $E=210\,\mathrm{GPa}$ y $\nu=0{,}30$ como propiedades elásticas de
referencia para acero, sujetas a la selección definitiva del material.

Las fuerzas normales de cada perforación y las componentes tangenciales se obtuvieron
de los estados correlacionados de la Etapa~3C/4A. Se analizaron, separadamente,
el instante de máxima fuerza ($F_z=14{,}071\,\mathrm{kN}$,
$T=88{,}924\,\mathrm{N\,m}$) y el instante de máximo momento por fuerzas normales
($M_f=153{,}802\,\mathrm{N\,m}$, con $F_z=12{,}864\,\mathrm{kN}$ y
$T=79{,}063\,\mathrm{N\,m}$). No se combinaron máximos correspondientes a instantes distintos.
Cada fuerza se distribuyó sobre una corona de material adyacente a su perforación
activa, con radios entre $3$ y $5{,}5\,\mathrm{mm}$. Esta distribución superficial
es una hipótesis aproximada de la transferencia de carga de la biomasa.

Se estudiaron dos condiciones de borde idealizadas: (i) la fijación completa de
la cara anular inferior entre diámetros $50$ y $58\,\mathrm{mm}$ y (ii) un apoyo
axial del mismo anillo combinado con restricción tangencial en un sector del
alojamiento central que representa provisionalmente la chaveta. Las condiciones
no reemplazan el contacto y la retención axial del montaje real.

\begin{table}[htbp]
\centering
\caption{Resultados preliminares FEM con malla de $90\,600$ tetraedros; unidades de tensión y desplazamiento.}
\label{tab:fem_matriz_4D}
\small
\begin{tabular}{llrr}
\hline
Caso & Restricción idealizada & $\sigma_{\mathrm{VM,max}}$ [MPa] & $|u_z|_{\max}$ [mm] \\
\hline
Máxima fuerza & Anillo fijo & 51{,}05 & 0{,}01182 \\
Máxima fuerza & Anillo axial + sector chaveta & 77{,}55 & 0{,}01355 \\
Máximo momento & Anillo fijo & 49{,}74 & 0{,}01094 \\
Máximo momento & Anillo axial + sector chaveta & 67{,}86 & 0{,}01265 \\
\hline
\end{tabular}
\end{table}

\begin{figure}[htbp]
\centering
\includegraphics[width=0.76\linewidth]{figuras/convergencia_desplazamiento.png}
\caption{Sensibilidad del desplazamiento vertical al refinamiento de malla, caso de máxima fuerza.}
\label{fig:conv_u_matriz}
\end{figure}

El desplazamiento máximo difiere menos de $1\%$ entre las dos mallas más finas,
mientras que las tensiones locales son sensibles al refinamiento y a las
restricciones en la región central. Además, en ambas hipótesis aparecen reacciones
negativas sobre sectores del apoyo inferior; este resultado cuestiona la
hipótesis de contacto bilateral sin retención axial y respalda incorporar una
fijación positiva para mantener el ensamble en servicio. No se interpretaron
estas reacciones como fuerzas de dimensionamiento directo de la tuerca.

El modelo no incorpora todavía los avellanados de $9\,\mathrm{mm}$, la ranura de
chaveta, sus concentraciones geométricas, el contacto unilateral con el hombro
ni la precarga de la retención. Por ello, los valores de la
Tabla~\ref{tab:fem_matriz_4D} permiten evaluar tendencias de rigidez y
sensibilidad, pero no establecen aún factores de seguridad definitivos.
Para finalizar la selección del material será necesario refinar estas regiones,
verificar convergencia de tensiones, documentar las propiedades mecánicas del
producto efectivamente suministrado y considerar desgaste, higiene y costo.
